\section{Pure owner covariance and moving endpoint readings}
\label{sec:covariant-owner}
We work in the full-owner primitive transversal $U_B$
from Theorem~\ref{thm:joint-transversal}, with
$B\subseteq M/E$ finite and $G$-invariant,
$N_B=\pi_E^{-1}(B)$, the thin group $G$, and
a chosen subgroup $\Gamma\le\Aut(B)$ acting purely on owner tags.
For the moving readings below, assume that $\Gamma$ contains the
owner action of every global thin translation in $G$; the action
on owners need not be faithful. The default choice is $\Aut(B)$,
but proper subgroups will be needed when a further owner quotient
must be preserved.
Its finite sorted strong exact families are
$\mathcal H_{st}$, with palettes $P_{st}=L_{st}\setminus F$.
The arguments below use only finite sorted primitive SAP,
strong typed CCO, and the pure owner action; they do not
require a root frame or a general conditional groupoid.

\begin{lemma}[Primitive loop saturation]\label{lem:covariant-loop-full}
In the full-owner transversal, for every nonempty loop exact
menu $W\in\mathcal H_{ss}$ one has
$W^{*n}=P_{ss}$ for all sufficiently large $n$.  Every compatible
nonempty exact menu absorbs a full palette on either side.
\end{lemma}
\begin{proof}
The universal primitive-loop lemma
\cref{lem:locality-positive} applies to the nonempty subset
$W\subseteq P_{ss}$ without requiring it to be original
exact.  Thus $W^{*n}=P_{ss}$ eventually.  The two-sided
full-palette absorption for every nonempty compatible menu,
including nonexact menus, is exactly
\cref{lem:locality-absorb}.  Both conclusions therefore
hold in the full-owner primitive transversal without another
proof of stabilization or a separate root-frame argument.
\end{proof}

Call a nonempty tagged menu \emph{proper} if it is not its full
pair palette.  For an original strong exact menu $C$, define
$\mathfrak F_\Gamma(C)$ to be the least finite family of
original exact tagged menus containing $C$ and closed under
converse, compatible left and right multiplication by \emph{all}
original exact menus, and the pure owner action $\Gamma$.
An absent multiplication factor is only an identity operation,
not an adjoined equality menu in the strong loop palette.

\begin{lemma}[No positive return through a proper menu]
\label{lem:covariant-no-return}
A directed closed walk in this generation graph through a
proper menu cannot use a multiplication move.  If two parallel
menus generate the same $\Gamma$-closed family, then they are
equal as actual raw-colour subsets.

More quantitatively, put
$m_0=\max_t|P_{tt}|\ge1$.
If a closed-word transformation $T$ has a genuine
nonempty positive outside multiplier, then
$T^{4|\Gamma|m_0}(D)=P_{st}$ for EVERY nonempty
parallel atomic menu $D\subseteq P_{st}$.
Here all compositions are STRONG and the exponent
counts iterations of $T$, not edges in an HKN reservoir.
\end{lemma}
\begin{proof}
Associate to each path its transformation of the distinguished
input menu.  By associativity of strong multiplication, by converse
reversing products, and by transporting factors under $\Gamma$,
every return to its initial pair of sorts has the normal form
\[
       T(X)=L*(\gamma X)^\varepsilon*R,
\]
where $\gamma\in\Gamma$, $\varepsilon$ denotes either the
original or the converse orientation, and each outside factor is
a nonempty original exact menu or is absent.  These factors have
the actual type labels forced by the path.  If the walk uses a
multiplication, at least one positive factor remains.  More
explicitly, each elementary multiplication contributes a
\emph{nonempty positive} typed profile to an outside word.
Pure owner relabelling merely transports the profile, while
converse reverses the order and exchanges the two outside words.
Associativity concatenates such words; by SAP their strong products
are nonempty.  Thus no normal-form rewrite can erase the last
positive outside factor or interpret it as an artificial identity.

Take a positive even multiple $m$ of the order of $\gamma$.
Upon composing $T$ with itself $m$ times, the central action
on $X$ becomes the identity, and the returning operation is
$T^m(X)=P*X*Q$ with loop multipliers at the original endpoints,
possibly absent, and with at least one nonempty loop multiplier.
At least one of these loop factors is present if the original
walk had a multiplication move: repetition, transport, converse and
positive concatenation preserve the existence of a positive outside
word.  These are genuine loop exact menus; all such products are
nonempty by SAP.  For an EXPLICIT uniform iteration bound, take
$q=2|\Gamma|$: this clears the single global pure
$\gamma$ and the possible converse in EVERY such return,
so $T^q(X)=P*X*Q$ for a pair of original-endpoint
loop multipliers, at least one of them nonempty positive.
Now iterate this transformation $k=2m_0$ more times.
By the uniform primitive loop-power bound
\cref{lem:locality-positive}, each PRESENT multiplier
$P^{*k}$ or $Q^{*k}$ is the full typed loop palette
at its endpoint sort.  Hence
\[
       T^{4|\Gamma|m_0}(D)
         = P^{*k}*D*Q^{*k}=P_{st}
\]
for every nonempty $D\subseteq P_{st}$, omitting
any absent multiplier as an IDENTITY ACTION only.
Full-palette absorption
\cref{lem:locality-absorb} proves the final equality.
If $T(D)=D$ for a PROPER D, this quantitative
conclusion is a contradiction.  Neither a synthetic
STRONG equality atom nor an invariant choice of HKN
packet origins entered the argument.

In particular a full palette cannot reach a proper one, by
absorption.  If parallel menus $A,B$ generate the same family,
there are paths $A\to B$ and $B\to A$.  If $A$ is proper,
their concatenation is a closed walk and contains no
multiplication.  Only pure type permutations and converse
remain; these fix the underlying raw atom names, so parallel
$A,B$ are equal.  The full case is immediate.
\end{proof}

\begin{remark}[Pure-tag covariance still needs symmetric RAW fine names]
\label{rem:covariant-converse-symmetry-firewall}
The final PARALLEL-menu step of
\cref{lem:covariant-no-return}
uses a stronger fact than
the condition that $\Gamma$
fixes fine-colour symbols:
in the ACTUAL owner transversal
all old fine atoms come from one
SYMMETRIC original binary structure.
Consequently taking typed converse
reverses the endpoint OWNER TAGS
but fixes the LITERAL underlying
old fine-colour name. Pure
$\Gamma$-relabelling likewise
fixes that name. A closed
sequence consisting solely of
these operations cannot change
a menu of parallel literal
fine-colour names.

This is FALSE in an arbitrary
ASYMMETRIC typed language,
even with SAP and CCO.
For a concrete negative
scope check, take two infinite
named sorts $U,V$, one
unconstrained symmetric
same-sort diversity atom
and two cross-sort DIRECTED
diversity atoms $a,b$ whose
typed converse interchanges
$a$ and $b$. Permit EVERY
consistent complete coloured
diagram (all typed atomic
triangles are legal).
The resulting random two-sort
age has free SAP and semantic
homomorphic CCO. Let
$\Gamma=C_2$ exchange $U,V$
while keeping the literal
directed atom names $a,b$.
On the cross-sort full menu
$D=\{a,b\}$, pure covariance
would give
\[
       h_{VU}(D)=h_{UV}(D),
\]
whereas converse compatibility
would give
\[
       h_{VU}(D)=\overline{h_{UV}(D)}.
\]
Neither $a$ nor $b$ is
self-converse, so no such
covariant singleton selector
exists. Equivalently the
parallel singleton menus
$\{a\}_{UV}$ and $\{b\}_{UV}$
generate the same
$\Gamma$-and-converse-closed
family WITHOUT any positive
multiplication; parallel-menu
separation is genuinely false.

Our ACTUAL symmetric-original
scope excludes this example.
One must NOT export the
present $\Gamma$-selector
argument to directed or
asymmetric original atoms
without an additional
converse-invariance criterion.
\end{remark}

\begin{theorem}[Pure-owner covariant joint hereditary selector]
\label{thm:covariant-HS}
The full strong exact menu families of $U_B$ admit one joint
hereditary selector $h$ preserving converse, all three PC
coordinates, and survivor heredity, and satisfying
\begin{equation}\label{eq:covariant-HS}
        h_{\gamma s,\gamma t}(D)=h_{st}(D)
       \qquad(\gamma\in\Gamma).
\end{equation}
\end{theorem}
\begin{proof}
The actual pure owner group $\Gamma$ permutes only
owner tags and fixes the literal fine-colour names.  It therefore
preserves the original exact-menu families, typed strong products,
converse and all legal triangle tables, as required by
\cref{lem:typed-pure-action-gate-engine}.  The
$\Gamma$-closed family in that lemma is precisely
$\mathfrak F_\Gamma(C)$ defined above.
Lemma~\ref{lem:covariant-no-return} supplies the required
no-return implication for parallel generated menus.
Consequently the common finite gate iteration produces a joint
hereditary singleton selector which preserves converse and all
three PC coordinates and whose gates commute with the pure
$\Gamma$-action.  This is exactly
\eqref{eq:covariant-HS}; no independent selection on different
owner pairs or averaging of selectors is involved.
\end{proof}

\section{Moving endpoint types and independent packet readings}
\label{sec:moving-readings}

For every $g\in G$, switching only the first endpoint of a fine
raw edge translates its owner $s$ to $gs$ and its colour $a$ to
$ga$, where $ga$ denotes the named thin translation action on
raw colours.  By the thin transport law, these operations send
legal triangles to legal triangles.  They preserve strong word
profiles, nonempty finite meets and the strong typed family:
\[
        D_{st}\longmapsto (gD)_{gs,t}.
\]
There is an analogous switch at the second endpoint.  Switching
both endpoints by the same $g$ fixes \emph{every raw fine colour}
and is the language action on types given by the diagonal
translation in $\Gamma$.  Define
\begin{equation}\label{eq:shifted-reading}
         h^g_{st}(D)=g\,h_{gs,t}(gD)
         \qquad(g\in G,\ D\in\mathcal H_{st}).
\end{equation}
It is crucial to change the endpoint type in evaluating $h$;
the unshifted expression $g\,h_{st}(gD)$ is generally ill-typed.

\begin{proposition}[Moving-type rectangular readings]
\label{prop:moving-readings}
The families $h^g$ are conservative, survivor-hereditary and
converse-compatible, are pure-owner $\Gamma$-covariant, and
obey the switching identity
\begin{equation}\label{eq:reading-switch}
    h^{u+v}_{vs,t}(vD)=v\,h^u_{st}(D).
\end{equation}
Moreover, if $D\in\mathcal H_{st}$, $A\in\mathcal H_{tu}$,
$P\in\mathcal H_{su}$ form a PC triple, then for
\emph{independently} chosen $g,k,\ell\in G$ their selections
form a legal fine-colour triangle:
\begin{equation}\label{eq:reading-rectangular}
   T\bigl(h^g_{st}(D),h^k_{tu}(A),h^\ell_{su}(P)\bigr).
\end{equation}
No distinctness of the sort labels is assumed.  The three
formal vertices of the triangle remain distinct even when their
owner labels coincide.
\end{proposition}
\begin{proof}
Conservativity and survivor heredity follow by transporting
$h$ across an endpoint translation, which bijectively preserves
inclusion of exact menus.  Since switching both endpoints by
$g$ fixes fine colours and the diagonal type action lies in
$\Gamma$, covariance~\eqref{eq:covariant-HS} gives
\[
 h_{gs,t}(gD)=h_{s,gt}(gD).
\]
Together with converse compatibility of $h$ this proves
converse compatibility of each $h^g$.  Any $\gamma\in\Gamma$
commutes with $G$ on the owners, so the same covariance formula
gives $\Gamma$-covariance of the readings.  The switching
identity is a direct calculation using exponent two:
\[
 h^{u+v}_{vs,t}(vD)
 =(u+v)h_{us,t}(uD)=v\,h^u_{st}(D).
\]

We first prove an auxiliary \emph{bilateral-support} statement:
if original exact menus $D_{st},A_{tu},P_{su}$ satisfy
\[
 D\subseteq P*\bar A,\qquad A\subseteq\bar D*P,
\]
then for every $g,k\in G$ one has
\begin{equation}\label{eq:reading-bilateral}
 P\cap\bigl(h^g(D)*h^k(A)\bigr)\ne\varnothing.
\end{equation}
Indeed $P'=P\cap(D*A)$ is nonempty exact.  For nonemptiness,
choose $d\in D$.  The inclusion $D\subseteq P*\bar A$ gives
$p\in P$ and $a\in A$ with $d\in p*\bar a$, hence
$p\in d*a\subseteq D*A$.  Nonempty exact intersections are exact
by Lemma~\ref{lem:typed-exact-products}.

We check the other two PC coordinates explicitly.  Strong SAP gives
the duplication inclusions
\[
 D\subseteq(D*A)*\bar A,\qquad
 A\subseteq\bar D*(D*A):
\]
for an arbitrary atom on either indicated side, attach an edge
with any atom of the other nonempty menu at a fresh third vertex.
Distributivity for the \emph{original exact} menus gives
\[
\begin{aligned}
 P'*\bar A
  &=(P*\bar A)\cap((D*A)*\bar A)\supseteq D,\\
 \bar D*P'
  &=(\bar D*P)\cap(\bar D*(D*A))\supseteq A.
\end{aligned}
\]
Together with $P'\subseteq D*A$, this is precisely
$\mathrm{PC}(D,A,P')$.  Switching its two outer endpoints by $g,k$
produces the PC triple
$(gD,kA,(g+k)P')$ of new sorts $(gs,t,ku)$.
The joint selector $h$ chooses a legal triangle for these menus.
Switching the outer endpoints back gives~\eqref{eq:reading-bilateral},
using the equality of the left and right shifted readings.

For a PC triple $(D,A,P)$ fix independently
$a=h^g(D)$, $b=h^k(A)$, $c=h^\ell(P)$.
Apply~\eqref{eq:reading-bilateral} to the Peirce rotation
with selected sides $(D,P)$ to obtain a triangle with \emph{these
particular} $a,c$ and some third colour in $A$; thus
$a\in\{c\}*\bar A$.  Similarly, applying it to $(A,P)$
gives $b\in\bar D*\{c\}$.
Consequently the exact nonempty traces
\[
 D'=D\cap(\{c\}*\bar A),\qquad
 A'=A\cap(\bar D*\{c\})
\]
contain respectively $a,b$.  Survivor heredity for the
\emph{individual readings} gives $h^g(D')=a$ and
$h^k(A')=b$.  Every colour in $D'$ has a partner in $A'$
forming a legal triangle with $c$, and conversely; hence
$D'\subseteq\{c\}*\overline{A'}$ and
$A'\subseteq\overline{D'}*\{c\}$.
Apply~\eqref{eq:reading-bilateral} to $(D',A',\{c\})$.
Its conclusion is that the selected $a,b,c$ themselves form a
legal triangle, proving~\eqref{eq:reading-rectangular}.
\end{proof}

\begin{remark}[Scope of the action and the reading]
The theorem uses only the pure $\Gamma$-action on the \emph{full}
owner set $B$.  It does not assert equivariant frames on the
previously chosen root representatives $S_0$.  There is no
invariant singleton hypothesis for a diagonal colour and no
independent choice of selectors for separate type pairs.
The independent shifts in~\eqref{eq:reading-rectangular} are
essential later for three distinct valuations to decode a legal
triangle.
\end{remark}
